\section{Vorstellung Medusa}

\begin{frame}{Medusa}
    \centering
    \textit{Multi‑dimensional Extrusion \& Dimensional Unified Slicing Algorithm}
    \vspace{0.4cm}

    \begin{columns}[T]
        \column{0.33\textwidth}
            \centering
            \faFileCode\\
            \textbf{Modell‑Darstellung}\\
            \small Import \& Visualisierung (STL/OBJ)
        \column{0.33\textwidth}
            \centering
            \faCogs\\
            \textbf{Slicing}\\
            \small Nichtplanare, parametrische Schichten
        \column{0.33\textwidth}
            \centering
            \faRobot\\
            \textbf{Pfad‑Generierung}\\
            \small 6‑Achs Toolpaths \& Ausgabeformate
    \end{columns}

    \vspace{0.35cm}
    \begin{block}{Kernaufgaben}
        \small
        \begin{itemize}
            \item Robustes Mesh‑Handling und Visual Debugging
            \item Modulare Slicing‑Pipeline (austauschbar)
            \item Integration mit KRONOS für Kinematik und Collision‑Checks
        \end{itemize}
    \end{block}
\end{frame}

    \begin{frame}{Medusa: Architektur \& Modulaufbau}
        \textbf{Layered Architektur:}
        \begin{itemize}
            \item \textbf{App}: Einstiegspunkt, Window/Main-Loop, Orchestrierung
            \item \textbf{UI (ImGui)}: Panels, File-Browser, Parameter/Settings
            \item \textbf{Graphics (OpenGL)}: Mesh, Shader, Renderer, Scene
            \item \textbf{Controls}: Input-Handling (Kamera/Model-Controller)
            \item \textbf{Slicing (geplant)}: Slicing-Pipeline \& 6-Achs Toolpaths
            \item \textbf{Core}: Logger, Camera, Math/Utilities (wiederverwendbar)
        \end{itemize}

        \vspace{0.35cm}
        \textbf{Kopplung/Prinzipien:}
        \begin{itemize}
            \item klare Modulgrenzen, erweiterbare Pipeline (Slicing austauschbar)
            \item Dependencies via CMake FetchContent (kein manuelles Setup)
        \end{itemize}
    \end{frame}

    \begin{frame}{Medusa: Architektur (Diagramm)}
        \centering
        \begin{tikzpicture}[
            font=\scriptsize,
            box/.style={draw=SecondaryColor, thick, rounded corners, align=center, minimum width=3.6cm},
            layer/.style={box, fill=SecondaryColor!8, minimum height=0.75cm},
            core/.style={box, fill=SecondaryColor!14, minimum height=0.75cm},
            ext/.style={box, fill=SecondaryColor!5, minimum height=0.75cm},
            mod/.style={box, fill=SecondaryColor!10, minimum height=0.6cm, minimum width=1.65cm},
            arrow/.style={->, thick, SecondaryColor!70}
        ]
            \node[layer] (app) {Application\\\texttt{src/app}};
            \node[mod, below left=0.6cm and 1.5cm of app.south] (controls) {Controls\\\texttt{src/controls}};
            \node[mod, below=0.6cm of app.south] (graphics) {Graphics\\\texttt{src/graphics}};
            \node[mod, below right=0.6cm and 1.5cm of app.south] (ui) {UI\\\texttt{src/ui}};

            \node[mod, below=0.25cm of graphics.south] (slicing) {Slicing (planned)\\\texttt{src/slicing}};
            \node[core, below=0.6cm of slicing.south, minimum width=6.8cm] (core) {Core\\Logger \textbar\ Camera \textbar\ Math\\\texttt{src/core}};
            \node[ext, below=0.6cm of core.south, minimum width=8.5cm] (ext) {External Libraries\\GLFW \textbar\ GLAD \textbar\ GLM \textbar\ ImGui \textbar\ Assimp \textbar\ spdlog};

            \draw[arrow] (app) -- (controls);
            \draw[arrow] (app) -- (graphics);
            \draw[arrow] (app) -- (ui);
            \draw[arrow] (graphics) -- (slicing);
            \draw[arrow] (controls) -- (core);
            \draw[arrow] (ui) -- (core);
            \draw[arrow] (slicing) -- (core);
            \draw[arrow] (core) -- (ext);
        \end{tikzpicture}
    \end{frame}
    
    \begin{frame}{Code Overview}
        \small
        \begin{columns}
            \column{0.55\textwidth}
                \begin{itemize}
                    \item \textbf{Total files:} 67
                    \item \textbf{Total lines of code:} 4\,286
                \end{itemize}

            \column{0.45\textwidth}
                \begin{tabular}{@{}l r@{}}
                \textbf{Language} & \textbf{LOC} \\\hline
                C++ & 2\,492 \\
                CMake & 786 \\
                C/C++ Header & 411 \\
                Markdown & 286 \\
                YAML & 257 \\
                Dockerfile & 54 \\\hline
                \textbf{SUM} & \textbf{4\,286}
                \end{tabular}
        \end{columns}
    \end{frame}
